% Part: applied-modal-logics
% Chapter: temporal-logic
% Section: possible-histories

\documentclass[../../../include/open-logic-section]{subfiles}

\begin{document}

\olfileid{aml}{tl}{poss}

\olsection{可能历史}

我们迄今使用的时态逻辑关系模型极具灵活性，因为我们无需对可达关系施加任何限制。这意味着时态模型可以在过去和未来分支，但我们或许想考虑一种更具“模态”色彩的分支理解，即将事件序列视为\emph{可能历史}。这未必需要改变我们的语言，不过我们也可以引入“普通的”模态算子 $\Box$ 和~$\Diamond$，还可以考虑引入认知可达关系以表示主体知识随时间的变化。

\begin{defn}
  时态语言的\emph{可能历史模型}是一个三元组 $\mModel{M} = \tuple{T, C, V}$，其中
  \begin{enumerate}
    \item $T$ 是一个非空集合，解释为时间中的状态。
    \item $C$ 是一个计算路径的集合，或称系统的\emph{可能历史}。换言之，$C$ 是由状态 $s_1$, $s_2$,~$s_3$, \dots 组成的序列~$\sigma$ 的集合，其中每个 $s_i \in T$。
    \item $V$ 是一个函数，它给每个命题变元~$p$ 指派一个时间点的集合~$V(p)$。
  \end{enumerate}
  为简化讨论，我们一般还假设，当一个历史在~$C$ 中时，它的所有后缀也在~$C$ 中。例如，如果 $s_1$, $s_2$,~$s_3$ 是~$C$ 中的一个序列，那么 $s_2$, $s_3$ 和~$s_3$ 也在~$C$ 中。此外，当两个状态 $s_i$ 和~$s_j$ 出现在序列~$\sigma$ 中时，若 $i < j$，则称 $s_i \prec_\sigma s_j$。当 $t \in V(p)$ 时，我们说 $p$~\emph{在~$t$ 处为真}。
\end{defn}

这里涉及的一个变化是，当我们在模型~$\mModel M$ 中评估公式在时间点~$t$ 的真值时，我们要相对于 $t$ 作为状态出现的某个历史~$\sigma$ 来进行。不过，我们无需改变命题变元或真值函数性联结词的语义。这些都与 \olref[sem]{defn:tmodels} 中的完全一样，因为它们都不会涉及~$\sigma$。然而，我们现在要相对于这些历史重新定义未来算子~$\Ftemp$，并引入~$\Diamond$ 算子。 

\begin{defn}\ollabel{defn:phmodels}
  \emph{公式~$!A$ 在~$\mModel M$ 中的 $t, \sigma$ 处为真}，记作 $\mSat{M}{!A}[t, \sigma]$，定义如下：
  \begin{enumerate}
  \item\ollabel{defn:sub:phmodels-f} \indcase{!A}{\Ftemp
    !B}{$\mSat{M}{\indfrm}[t, \sigma]$ 当且仅当 $\mSat{M}{!B}[t', \sigma]$，存在某个 $t' \in T$ 满足 $t \prec_\sigma t'$。}
  \item\ollabel{defn:sub:phmodels-diamond} \indcase{!A}{\Diamond
    !B}{$\mSat{M}{\indfrm}[t, \sigma]$ 当且仅当 $\mSat{M}{!B}[t, \sigma']$，存在某个 $\sigma' \in C$ 使得 $t$ 出现在 $\sigma'$ 中。}
  \end{enumerate} 
\end{defn}

其他时态和模态算子可类似定义。不过，我们现在能够表示结合了时态与模态的断言。例如，我们可以用公式 $\lnot \Ftemp p \land \Diamond \Ftemp p$ 来符号化“$p$~将不发生，但它本可能发生”。这将在满足如下条件的点和历史中成立：在该点和该历史中，$p$ 不会在后继状态中成真，但存在另一条历史，在其中 $p$ 将会成真。 

\end{document}